% indicator to skip certain packages when converting to HTML5 / ePub
\newif\ifnotlatexml
\ifdefined\latexml
  \notlatexmlfalse
\else
  \notlatexmltrue
\fi

% set page margins
\usepackage[left=1.25in, right=1.25in, top=1in, bottom=1.15in]{geometry}

% configure fonts for PDF exports
\ifnotlatexml
\usepackage{fontspec}
\setmainfont{Source Serif 4}[
    UprightFont = * Regular,
    ItalicFont = * Italic,
    BoldFont = * Semibold,
    BoldItalicFont = * Semibold Italic
]
\setsansfont{Source Sans 3}[
    UprightFont = * Regular,
    ItalicFont = * Italic,
    BoldFont = * Semibold,
    BoldItalicFont = * Semibold Italic
]
\setmonofont{Source Code Pro}[
  UprightFont = * Regular,
  ItalicFont = * Italic,
  BoldFont = * Semibold,
  BoldItalicFont = * Semibold Italic
]
\fi

% micro-typography: character protrusion and font expansion. Noticeably
% improves justified text by reducing loose lines and rivers.
\ifnotlatexml
\usepackage{microtype}
\fi

% format text spacing
\usepackage{setspace}
\setstretch{1.1}
\setlength{\parindent}{0pt}
% rubber, not rigid: gives the page-breaking algorithm room to flex the gap
% between paragraphs rather than stranding a lone line at a page boundary
\setlength{\parskip}{6pt plus 2pt minus 1pt}

% discourage widows (lone line at the top of a page) and orphans (lone line at
% the bottom). If pages start coming out noticeably short, lower these to ~4000.
\widowpenalty=10000
\clubpenalty=10000

% prevent paragraph skip from adding extra space immediately after headings
\makeatletter
\g@addto@macro\@afterheading{\vspace{-\parskip}}
\makeatother

% conditional logic --- loaded early because the colour block below needs
% \ifdefempty to treat the secondary brand colour as optional
\usepackage{etoolbox}

% custom colors
\usepackage{xcolor}

% define colorblind-friendly palette
\definecolor{cb1}{HTML}{56B4E9}
\definecolor{cb2}{HTML}{E69F00}
\definecolor{cb3}{HTML}{009E73}
\definecolor{cb4}{HTML}{0072B2}
\definecolor{cb5}{HTML}{D55E00}
\definecolor{cb6}{HTML}{CC79A7}
\definecolor{cb7}{HTML}{999999}
\definecolor{cb8}{HTML}{F0E442}

\definecolor{schoolcolorprimary}{HTML}{\brandhexprimary}

% Two brand colours, one job each:
%   schoolcolorprimary   -> section heading text
%   schoolcolorsecondary -> the rule beneath section headings
%
% The secondary is OPTIONAL. Institutions built on a single mark may leave
% \brandhexsecondary empty, in which case it aliases the primary and the rule
% matches the heading text. Guarding this matters: \definecolor{...}{HTML}{}
% on an empty value is a compile error, not a silent no-op.
\ifdefempty{\brandhexsecondary}{%
  \colorlet{schoolcolorsecondary}{schoolcolorprimary}%
}{%
  \definecolor{schoolcolorsecondary}{HTML}{\brandhexsecondary}%
}

% Format section headings.
%
% Headings are set in the sans family (Source Sans 3) against the serif body
% (Source Serif 4). The family shift carries most of the hierarchy, which lets
% the size ramp stay modest at a 10pt base: 14.4 / 12 / 10pt.
%
% Note \bfseries maps to Source Sans 3 *Semibold*, not a true Bold, so these
% read as authoritative rather than heavy.
%
% The label argument is {} because every heading in this system is starred
% (\section*). Restore {\thesection.} if numbered headings are ever wanted.
\usepackage{titlesec}
\titleformat{\section}
  {\Large\sffamily\bfseries\color{schoolcolorprimary}}
  {}
  {0.5em}
  {}
  [{\color{schoolcolorsecondary}\titlerule[0.4pt]}]

\titlespacing*{\section}
  {0pt}  % left
  {3pt} % before
  {6pt}  % after

% format subsection headings
\titleformat{\subsection}
  {\large\sffamily\bfseries}
  {}
  {0.5em}
  {}

\titlespacing*{\subsection}
  {0pt}  % left
  {2pt}  % before
  {2pt}  % after

% Subsubsections sit at body size, so the serif/sans contrast alone has to
% carry the level --- which it does more reliably than the previous serif
% italic, which disappeared into the surrounding text at 10pt.
\titleformat{\subsubsection}
  {\normalsize\sffamily\bfseries}
  {}
  {0.5em}
  {}

\titlespacing*{\subsubsection}
  {0pt}  % left
  {2pt}  % before
  {0pt}  % after

% avoid "widowed" headings / subheadings
\usepackage{needspace}

\let\oldsection\section
\renewcommand{\section}{\Needspace{5\baselineskip}\oldsection}

\let\oldsubsection\subsection
\renewcommand{\subsection}{\Needspace{3\baselineskip}\oldsubsection}

% Adjust line spread and font family inside "pull quotes"
\usepackage{csquotes}
\newenvironment{pullquote}[1][]%
  {\begin{displayquote}[#1]\sffamily}%
  {\end{displayquote}}

% write math and define operators
\usepackage{mathtools}
\ifnotlatexml
\usepackage{unicode-math}
\setmathfont{NewComputerModernMath}
\setmathfont{NewComputerModernMath}[range=\mathbb]
\newcommand{\indicator}{\mathbb{1}}
\else
\usepackage{amsfonts, amssymb, dsfont}
\newcommand{\indicator}{\mathds{1}}
\fi

% images and figures
\usepackage{graphicx, subcaption}

% links
\usepackage{hyperref}
% Link colour has to satisfy two requirements that pull in opposite directions:
%
%   vs. the white background : >= 4.5:1 (WCAG AA, body-size text) -> wants dark
%   vs. the black body text  : >= 3.0:1 (WCAG 1.4.1, which applies because
%                              colour is the only cue distinguishing a link
%                              here --- colorlinks adds no underline) -> wants light
%
% Together these confine luminance to roughly [0.100, 0.183]. 25708A sits mid-band
% at 5.58:1 against white and 3.76:1 against black, clearing both with real margin.
%
% For reference: 2A7F9E passed white by only 0.04; 1F6478 passed black by only
% 0.15; and 1A5A6B, which would have reached AAA against white, FAILS against
% black at 2.72:1 --- chasing AAA on the background breaks link legibility.
\definecolor{url}{HTML}{25708A}
\hypersetup{
  colorlinks = true,
  linkcolor = black,
  filecolor = black,
  urlcolor = url,
  citecolor = url
}

% tables
\usepackage{booktabs, longtable, threeparttable}
\usepackage{lscape} % landscape tables
\usepackage{array}
\newcolumntype{P}[1]{>{\centering\arraybackslash}p{#1}}

% Aligned label/value blocks (Instructor Information, Class Meeting Times,
% Course Details, ...). Each of these used its own tabular{@{}ll}, so every
% block sized its label column independently and the value columns started at
% a different horizontal position down the page. A shared fixed label width
% lines them all up.
%
% \deflabelwidth is sized for the longest label in the system, "Equivalent
% Courses:". Shorten it if a given institution's labels are all brief.
%
% The value column is a wrapping p-column spanning the rest of the measure, so
% long values (e.g. the email line with its subject-line note) reflow instead
% of running into the right margin.
\newlength{\deflabelwidth}
\setlength{\deflabelwidth}{1.15in}
\newenvironment{deftable}{%
  \begin{tabular}{@{}%
    >{\raggedright\arraybackslash}p{\deflabelwidth}%
    >{\raggedright\arraybackslash}p{\dimexpr\linewidth-\deflabelwidth-2\tabcolsep\relax}%
    @{}}%
}{%
  \end{tabular}%
}

% Keep a lead-in sentence with the centred block it introduces.
%
% \widowpenalty and \clubpenalty cannot do this: they only govern breaks
% *within* a paragraph, so a one-line lead-in ("...using the following scale:")
% has no internal break to protect and gets stranded at the foot of a page
% while its table moves to the next.
%
% Forbidding the break immediately before every centred block forces TeX to
% break at the previous legal point instead --- i.e. before the lead-in --- so
% sentence and table travel together. No markup needed in the document source.
%
% If a centred block ever grows taller than a page, soften [4] to [3], which
% strongly discourages rather than forbids.
%
% Display math needs no equivalent: TeX's \predisplaypenalty already defaults
% to 10000, so a break just before \[ ... \] is forbidden out of the box.
\BeforeBeginEnvironment{center}{\par\nopagebreak[4]}

% nice enumerate environments in tagged PDFs
\usepackage{enumitem}
\usepackage{multicol}

% configuring global itemize doesn't work when tagging,
% so make "custom" environment first
\newlist{ul}{itemize}{3}

\setlist[ul,1]{
  label=\textbullet,
  topsep=2pt,
  itemsep=1pt,
  parsep=0pt,
  partopsep=0pt,
  leftmargin=2.5em
}

\setlist[ul,2]{
  label=\textendash,
  topsep=1pt,
  itemsep=1pt,
  parsep=0pt,
  partopsep=0pt,
  leftmargin=2.5em
}

\setlist[ul,3]{
  label=\textasteriskcentered,
  topsep=1pt,
  itemsep=0pt,
  parsep=0pt,
  partopsep=0pt,
  leftmargin=2.5em
}

% configuring global enumerate doesn't work when tagging,
% so make "custom" environment first
\newlist{ol}{enumerate}{3}

\setlist[ol,1]{
  label=\arabic*.,
  topsep=3pt,
  itemsep=2pt,
  parsep=0pt,
  partopsep=0pt,
  leftmargin=2.5em
}

\setlist[ol,2]{
  label=(\alph*),
  topsep=2pt,
  itemsep=1pt,
  parsep=0pt,
  partopsep=0pt,
  leftmargin=2.5em
}

\setlist[ol,3]{
  label=\roman*.,
  topsep=1pt,
  itemsep=0pt,
  parsep=0pt,
  partopsep=0pt,
  leftmargin=2.5em
}


